\begin{proposition}[Arity-one counterterm and obstruction along $\Delta$]
\label{prop:c9-structure}
\begin{itemize}
\item (i) At one loop, $I^{CT}_{\mathfrak{psl}(4|4)}[\varepsilon]^{(1)}(\Delta)$ is the singular part of the composite cycles $\tfrac\hbar2\,\partial_{2(P\star v.P\star\otimes1)K_0}W(P(\varepsilon,\infty),I)_{\mathrm{tree}}$, $P=P(\varepsilon,\infty)$, in which the composite kernel closes two legs of a tree of $I$; it is $\log\varepsilon$ times a local functional, with no coefficient of $\varepsilon^{-m}$, $m\ge1$. Every other one-loop graph with one vertex $S^{(1)}(\Delta)$ has a limit as $\varepsilon\to0$ at generic momenta.
\item (ii) (Local formula.) Assume that $\hbar F_1^{(0)}$ solves the arity-zero part of $s(\hbar F_1)=\hbar\,\mathrm{Ob}$ in the sectors with at most $n$ legs. A defect cycle along $\Delta$ is a cycle all of whose vertices lie on it, with exactly one special edge, exactly one vertex $S^{(1)}(\Delta)$ and the other vertices from $I$, the vertex $S^{(1)}(\Delta)$ placed on a line, its term $(v.A|_V,A^*|_V)$ with both legs on edges $P$, which is the composite kernel, or at the special edge, with one leg on it. Each defect cycle $\mathcal C$ with external legs at generic momenta is $\sum_{m\ge1}\varepsilon^{-m}a_m(\mathcal C)+a_0(\mathcal C)+o(1)$ with local coefficients and no term in $\log\varepsilon$, and in the sectors with $n$ legs
\begin{equation}
\hbar\,\mathrm{Ob}^{(1)}(\Delta)\big|_{n\ \text{legs}}=\sum_{\mathcal C\ \text{with}\ n\ \text{legs}}a_0(\mathcal C),\qquad \sum_{\mathcal C\ \text{with two legs}}a_1(\mathcal C)=\tfrac94\,\text{units}\cdot\Big\{\int\mathrm{dvol}\,(A,A),S^{(1)}(\Delta)\Big\},
\label{eq:c9-local}
\end{equation}
the other pole sums vanishing. This is the definition of the numbers of the obstruction.
\item (iii) (Second evaluation.) Assume moreover that the $\varepsilon\to0$ limit of $Q$ of the composite cycles minus $(Q+\{I,-\})I^{CT}_{\mathfrak{psl}(4|4)}[\varepsilon]^{(1)}(\Delta)$ is the tree-level flow of a local functional, the composite part. Then, in the sectors with at most $n$ legs,
\begin{equation}
\hbar\,\mathrm{Ob}^{(1)}(\Delta)=4\big(\text{coefficient of }\log\varepsilon\text{ in }I^{CT}[\varepsilon]\big)+(\text{composite part})=(Q+\{I,-\})\Big(4\,\frac{\hbar\operatorname{tr}_{\mathfrak g}(\mathrm{ad}_X\mathrm{ad}_Y)}{16\pi^2(X,Y)}\Big(2\int(A,A^*)-4\int(\lambda^*,\lambda)-\int(c^*,c)\Big)\Big)+(\text{composite part}).
\label{eq:c9-Ob}
\end{equation}
\end{itemize}
\end{proposition}

\begin{proof}
(i) By \eqref{eq:c9-D1a} at $T=\infty$, the one-loop graphs with one vertex $S^{(1)}(\Delta)$ are $(P\star v.)$ applied to the one-loop part of $W(P,J)$, the composite cycles, and the graphs joining the tree part and the one-loop part of $W(P,J)$ by the composite kernel. The one-loop part of $W(P,J)$ has a limit at generic momenta (Proposition \ref{prop:c4-limit}), and the first and third classes attach to its legs an edge $P$ or the composite kernel at generic momenta, so they have limits; the counterterm, the singular part of the graphs (Definition \ref{def:c8-quant}), is that of the composite cycles. A coefficient of $\varepsilon^{-m}$ of $I^{CT}_{\mathfrak{psl}(4|4)}[\varepsilon]^{(1)}(\Delta)$ is a colour-singlet local functional of degree $-1$ and Weyl weight $4-2m$, the position at weight $-1$ and at most linear in it; a colour singlet of degree $-1$ contains an antifield of weight at least $\tfrac52$ and a second adjoint field of weight at least $1$, or $c^*$ and $c$, so its total weight is at least $\tfrac52$, above $4-2m\le2$.

(ii) By Lemma \ref{lem:c8-transfer} the obstruction is read from the $\varepsilon\to0$ limit of the connected graphs with one vertex from the defect vertex of Definition \ref{def:c9-data}, the other vertices from $I-I^{CT}[\varepsilon]-\hbar F_1^{(0)}$ and $S^{(1)}(\Delta)$, edges $P(\varepsilon,\infty)$. At order $\hbar$ a graph has one cycle and vertices of order $\hbar^0$, or no cycle and one vertex of order $\hbar$ (Proposition \ref{prop:c3-graphs}). The vertices of order $\hbar^0$ of the defect vertex are special-edge parts of brackets, so the special edge lies on the cycle or on a tree. The graphs are: the defect cycles; the reducible graphs, a defect cycle with fewer legs or a one-loop subgraph of the arity-zero defect joined by edges $P$ to trees of $I$ and $S^{(1)}(\Delta)$; the graphs whose special edge is a bridge; the trees with one counterterm vertex; and the constant $\hbar\Delta_\varepsilon S^{(1)}(\Delta)$. A bridge at a generic momentum $p$ separates a one-loop subgraph $G$; with the special-edge part of the counterterm bracket across the same edge, its counterterm partner, the sum is $(e^{-\varepsilon|p|^2}-1)(G-\mathrm{Sing}\,G)$, $e^{-\varepsilon|p|^2}-1$ times a convergent graph, which tends to $0$; the partners carrying $\log\varepsilon$ give $(e^{-\varepsilon|p|^2}-1)\log\varepsilon\to0$. By induction on the number of legs the reducible graphs tend to the trees of the components of the obstruction with fewer legs, the arity-zero two-leg subgraphs being Theorem \ref{thm:c4-obstruction}. A defect cycle is a sum of simplex integrals over the region where the Schwinger coordinate of the special edge is the smallest, the domains with the cutoff shapes $1/\min_{f\ne\sigma}\alpha_f$ and $1/\min_f\alpha_f$ differing there, so the terms in $\log\varepsilon$ cancel (Lemma \ref{lem:c10-special} and the argument of Theorem \ref{thm:c11-simplex}, which apply to cycles of vertices of $I$ with the composite kernel of order $|p|^{-3}$, Proposition \ref{prop:c9-composite} (iv)). The left side has a limit and the defect cycles carry no $\log\varepsilon$, so the terms in $\log\varepsilon$ of the counterterm trees cancel among themselves, their special-edge parts tending to $0$; the only pole of the counterterm trees is that of $-\{I^{CT}[\varepsilon]^{(0)},S^{(1)}(\Delta)\}_\varepsilon$, with $I^{CT}[\varepsilon]^{(0)}$ the counterterm of Theorem \ref{thm:c4-quadratic}, by (i). The cancellation of the poles gives the pole identity, and the tree-level flow of the part with $n$ legs of the obstruction is $\sum a_0(\mathcal C)$ plus the trees of the components with fewer legs; the injectivity of the tree-level flow (Theorem \ref{thm:c8-L} (i)) gives \eqref{eq:c9-local}. The vertex $\hbar F_1^{(0)}$ adds $-\{\hbar F_1^{(0)},S^{(1)}(\Delta)\}_\varepsilon$, whose limit $\{S^{(1)}(\Delta),\hbar F_1^{(0)}\}$ vanishes (Lemma \ref{lem:c9-SDelta} (i)).

(iii) Apply \eqref{eq:c9-D1b} at $T=\infty$, where $\{-,-\}_\infty=0$. At order $\hbar$: $\{S^{(1)}(\Delta),W\}_\varepsilon$ differs from $\{S^{(1)}(\Delta),W\}$ by a special edge on a leg, carrying $e^{-\varepsilon|p|^2}-1$ at generic momenta, and $\{S^{(1)}(\Delta),W\}=v.W$ tends by \eqref{eq:c9-ward}, as in the proof of Corollary \ref{cor:c9-CS}, to the tree-level flow of $4$ times the coefficient of $\log\varepsilon$ in $I^{CT}[\varepsilon]$. $QW$ is the flow of the arity-zero defect (Proposition \ref{prop:c4-defectflow}); its tree part carries a special edge on a tree and tends to $0$, and its one-loop part tends to $0$ by the arity-zero equation, so $(P\star v.)(QW)$ and the graphs joining $QW$ and $W$ tend to $0$, while $-\tfrac\hbar2\partial(QW_{\mathrm{tree}})=Q\big(\tfrac\hbar2\partial W_{\mathrm{tree}}\big)$ by Proposition \ref{prop:c9-composite} (iii). $Q$ of the linearised flow of $I^{CT}_{\mathfrak{psl}(4|4)}[\varepsilon]^{(1)}(\Delta)$ tends to the tree-level flow of its image under $Q+\{I,-\}$, as in Lemma \ref{lem:c6-bvexact}. Both (ii) and (iii) compute the tree-level flow of the same local functional, which is unique (Theorem \ref{thm:c8-L} (i)). The second form is \eqref{eq:c6-fieldren}.
\end{proof}
